\newpage
\section*{ЗАКЛЮЧЕНИЕ}
\addcontentsline{toc}{section}{ЗАКЛЮЧЕНИЕ}

В результате выполнения курсовой работы была разработана программная система для сжатия изображений с потерями в формате JPEG. Реализованное программное обеспечение позволяет выполнять сжатие цифровых изображений с заданной степенью качества, сохранять результаты обработки, а также анализировать эффективность сжатия по размеру выходных файлов.

В ходе работы был проведён анализ предметной области, связанной с методами и алгоритмами сжатия изображений. Особое внимание было уделено стандарту JPEG, основанному на использовании дискретного косинусного преобразования, квантования коэффициентов и субдискретизации цветовых компонент. Рассмотренные методы позволили сформировать обоснованные требования к функциональности разрабатываемой программы.

Разработанная система реализует полный цикл JPEG-сжатия, включающий:
\begin{itemize}
	\item преобразование цветового пространства RGB в YCbCr;
	\item субдискретизацию цветоразностных компонент по схеме 4:2:0;
	\item разбиение изображения на блоки размером $8 \times 8$;
	\item применение дискретного косинусного преобразования;
	\item квантование коэффициентов с учётом заданного уровня качества;
	\item восстановление изображения и сохранение результата в формате JPEG.
\end{itemize}

Для взаимодействия с пользователем был реализован графический интерфейс, обеспечивающий загрузку одиночных изображений и папок с файлами, выбор степени сжатия, отображение исходного и итогового размеров файлов, а также визуализацию процесса выполнения с помощью индикатора прогресса. Дополнительно предусмотрена возможность построения графика зависимости размера изображения от степени сжатия, что позволяет наглядно оценить влияние параметра качества на эффективность алгоритма.

В ходе разработки было проведено тестирование программы на различных изображениях и значениях качества. Результаты подтвердили корректность работы реализованного алгоритма JPEG-сжатия и соответствие получаемых результатов теоретическим ожиданиям.

Основные результаты курсовой работы:
\begin{enumerate}
	\item Проведён анализ методов сжатия изображений с потерями и стандарта JPEG.
	\item Спроектирована архитектура программной системы для JPEG-сжатия.
	\item Реализован алгоритм сжатия изображений с использованием дискретного косинусного преобразования и квантования.
	\item Разработан графический интерфейс для управления процессом сжатия и анализа результатов.
	\item Выполнено тестирование программы и оценена эффективность сжатия.
\end{enumerate}

Все задачи, поставленные в начале курсовой работы, были успешно решены. Таким образом, цель курсовой работы — разработка программы для сжатия изображений с потерями в формате JPEG с возможностью анализа степени сжатия — была достигнута.